\subsection{Cyclotomic bounds in every dimension and a character lower bound}

\begin{lemma}[B-RD14: full torsion forces the cyclotomic extension]
\label{lem:B-RD14}
\status{Proved.}
Let $A/K$ be an abelian variety of positive dimension over a
characteristic-zero field. If $A[p^n]$ is rational over $K$, then
$\mu_{p^n}\subset K$. A principal polarization, or a polarization of
degree prime to $p$, is not needed for this implication.
\end{lemma}
\begin{proof}
Choose any polarization defined over $K$. Its Weil pairing gives a
nonzero $G_K$-equivariant integral alternating form
\[
 b:T_pA\times T_pA\longrightarrow\mathbf Z_p(1).
\]
Divide all its coefficients by their largest common power of $p$.
The resulting form is still integral and equivariant, and at least
one coefficient is a unit. Since $G_K$ acts trivially on $T_pA/p^n$,
equivariance yields
$(\chi_p(\sigma)-1)b(x,y)\equiv0\pmod {p^n}$ for every $x,y$.
Choose a pair with unit value; then
$\chi_p(\sigma)\equiv1\pmod {p^n}$ for every $\sigma\in G_K$.
This is precisely the claimed cyclotomic containment.
\end{proof}

\begin{proposition}[B-RD15: the unconditional reading of the least-level question]
\label{prop:B-RD15}
\status{Proved; the final implication is explicitly vacuous.}
Fix $p\geq3$, a positive multiple $e$ of $p-1$, and $g\geq1$, and put
$s=v_p(e/(p-1))$. Among \emph{all} abelian varieties of dimension $g$
over all complete discretely valued fields with perfect residue field
and absolute ramification $e$, the least positive $n$ for which full
rational $p^n$-torsion implies good reduction is $s+2$.
At this level full rational $p^n$-torsion cannot occur at all.
Consequently a substantive good-reduction threshold must specify the
potentially good class, as in Corollary~\ref{cor:B-RD13}.
\end{proposition}
\begin{proof}
Lemma~\ref{lem:B-RD14} and the ramification degree
$[\mathbf Q_p(\mu_{p^n}):\mathbf Q_p]=p^{n-1}(p-1)$ imply
$p^{n-1}(p-1)\mid e$. Thus full rational $p^n$-torsion is impossible
for $n\geq s+2$.

Conversely, for each $1\leq n\leq s+1$, choose a complete field $K$
with residue field $\overline{\mathbf F}_p$, absolute ramification $e$,
and $\mu_{p^n}\subset K$. Such a field is obtained by taking the
cyclotomic extension of $\operatorname{Frac}W(\overline{\mathbf F}_p)$
and a further totally ramified extension of the required degree.
By \cite[Theorem~1, p.6 of the author version]{Tate1995Uniformization},
a split Tate curve with parameter $q=a^{p^n}$, where $a$ is a
uniformizer, has all its $p^n$-torsion represented by
$\zeta_{p^n}^i a^j$ in $K^\times/q^{\mathbf Z}$, and has split
multiplicative reduction. For dimension $g$, take its product with
$g-1$ canonical ordinary elliptic lifts. Their $p$-divisible groups
are $\mu_{p^\infty}\oplus\mathbf Q_p/\mathbf Z_p$, so they too have
full rational $p^n$-torsion. The product still has bad reduction.
This proves the asserted least level and explains its vacuity.
\end{proof}

\begin{proposition}[B-RD16: necessary lower bounds for character determination]
\label{prop:B-RD16}
\status{Proved.}
Let $h\geq2$. For $p\geq5$ and absolute ramification $e$, let $N$
be the largest bad elliptic level from Corollary~\ref{cor:B-RD13}.
If $N\geq1$, any universal least index $\nu$ in B.4 must satisfy
$\nu\geq N$. More precisely, for each $1\leq n\leq N$ there are two
height-$h$ BT groups satisfying the manuscript's monodromy condition,
with isomorphic $p^n$-torsion, whose characters on inertia elements of
order prime to $p$ differ. For $p=3$, the same lower bound holds using
the tame value $m_{\max}=4$ in that corollary.
\end{proposition}
\begin{proof}
Take the elliptic curve $E/K$ supplied by Theorem~\ref{thm:B-RD12}
at level $n$, and a canonical ordinary elliptic lift $E'/K$.
Lemma~\ref{lem:B-RD14} gives $\mu_{p^n}\subset K$, so both
$E[p^n]$ and $E'[p^n]$ are constant of rank $p^{2n}$.
Their inertia characters differ: the ordinary good curve has trace
two, whereas the nontrivial tame automorphism of the supersingular
good reduction has trace different from two. Append the same constant
etale BT group of height $h-2$ to both. Their characters each increase
by $h-2$, preserving the difference, and their truncated groups remain
isomorphic. These groups are potentially good, hence satisfy the
stated monodromy condition without using any extension argument.
Finally the manuscript's $M(\nu)$ is $M[p^{\nu+1}]$; failure at
$n=N$ gives $\nu\geq N$.
\end{proof}
